\documentclass{article}

\usepackage[english]{babel}

% Set page size and margins
% Replace `letterpaper' with `a4paper' for UK/EU standard size
\usepackage[letterpaper,top=2cm,bottom=2cm,left=1cm,right=1cm,marginparwidth=1.75cm]{geometry}

% Useful packages
\usepackage{amsmath}
\setcounter{MaxMatrixCols}{15}
\usepackage{graphicx}
\usepackage[colorlinks=true, allcolors=blue]{hyperref}
\usepackage{multicol}
\usepackage{fancyhdr}
\usepackage{caption}
\usepackage{braket}
\usepackage{mathtools}
\usepackage{tikz}

\title{8,3,2 Code}

\begin{document}

\section{Construction of the Code}
From this paper: https://doi.org/10.1103/PhysRevA.109.062438, the code is constructed on the 8 vertices of a cube, with the $Z$ stabilizers corresponding to faces of the cube, and the $X$ stabilizer is the global parity check using all vertices. The stabilizer is group is
\[S = \begin{pmatrix}
    Z & Z & Z & Z & I & I & I & I\\
    I & I & I & I & Z & Z & Z & Z\\
    Z & Z & I & I & Z & Z & I & I\\
    Z & I & Z & I & Z & I & Z & I\\
    X & X & X & X & X & X & X & X\\
\end{pmatrix}\]
According to the paper, the logical operators also have geometric interpretations. Logical $X$ operators are associated with faces of the cube, and logical $Z$ operators are associated with edges of the cube.

\section{Logical operators}
Given the stabilizers, I searched for logical operators and found 14 $X$ logical operators, each of weight 4, corresponding to the 6 faces of the cube, 6 planar diagonals normal to face diagonals, and 2 regular tetrahedra. I found 123 $Z$ logical operators, which includes the set of all even weight $Z$ operators, excluding the 4 $Z$ stabilizers and the empty operator:
\[123 = \sum_{k=0}^4 \binom{8}{2k} - 4 - 1= {8\choose0} - 1 + {8\choose2} + {8\choose4} - 4 + {8\choose6} + {8\choose8}\]
The lowest weight operators of these are the weight-2 operators which can represent edges of the cube. There isn't an apparent nice geometric interpretation for the $Z$ logical operators, it includes every possible combination that commute with the $X$ global parity check.

\section{Study Idea}
\subsection{Idea 1}
Start by gauging $Z$ logical operators, exhausting the gauge fixing space. Then, for each gauge fixing, measure all physical qubits in $Z$, and use the error rate on each to estimate the error rate of possible $Z$ logical operators.


\section{Exhaustive Logical Operator Promotion}

\section{Algebraic Model of Logical Operator Promotion}
I also created an algebraic model which assigns an equal probability of each error type to occur on each physical qubit. Then, assuming syndrome extraction and final readout is error free, I calculate the error rate for all possible $Z$ logical operators, I'm focusing on just $Z$ for now.

For example, consider the logical operator $ZZIIIIII$ ($Z_0Z_1$). This logical operator will error if an $X$ error occurs on qubits 0 xor 1. However, in the default code most full readouts in which this logical operator errors have nonzero syndromes, that is, the error is correctly detected by the stabilizers, I'll call these ``sneaky'' readouts. There are only a handful of sneaky readouts, some of which cause a logical error on this operator. 

In the default code, there are 14 sneaky readouts, these are possible final full readouts where physical qubits have flipped but cannot be detected by the ($Z$) stabilizers:
\[\begin{bmatrix}
    0 & 0 & 0 & 0 & 1 & 1 & 1 & 1\\
    0 & 0 & 1 & 1 & 0 & 0 & 1 & 1\\
    0 & 0 & 1 & 1 & 1 & 1 & 0 & 0\\
    1 & 1 & 0 & 0 & 0 & 0 & 1 & 1\\
    1 & 1 & 0 & 0 & 1 & 1 & 0 & 0\\
    1 & 1 & 1 & 1 & 0 & 0 & 0 & 0\\
    0 & 1 & 0 & 1 & 0 & 1 & 0 & 1\\
    1 & 0 & 1 & 0 & 1 & 0 & 1 & 0\\
    0 & 1 & 0 & 1 & 1 & 0 & 1 & 0\\
    0 & 1 & 1 & 0 & 0 & 1 & 1 & 0\\
    0 & 1 & 1 & 0 & 1 & 0 & 0 & 1\\
    1 & 0 & 1 & 0 & 1 & 0 & 1 & 1\\
    1 & 0 & 0 & 1 & 1 & 0 & 0 & 1\\
    1 & 0 & 1 & 0 & 0 & 1 & 0 & 1\\
\end{bmatrix}\]

Of these 14 sneaky readouts, 8 satisfy $q_0 \oplus q_1 = 1$. The probability of any of these happening is \[8p^4(1-p)^4\] which is consistent with the algebraic model when ignoring $Z$ errors detected by the $X$ stabilizer. Since half of all possible error 
\end{document}